% !TeX root = ../main.tex
% Add the above to each chapter to make compiling the PDF easier in some editors.

\chapter{Appendix}\label{chapter:appendix}

\section{Dataset Sources}
\label{appendix:datasets}
The datasets used in this thesis are sourced from various platforms, including Roboflow and Kaggle. The Table~\ref{tab:dataset_summary}
summarizes the datasets used for training, validation, and testing.
\begin{table}[htbp]
\centering
\small
%\renewcommand{\arraystretch}{2}
\resizebox{\textwidth}{!}{
\begin{tabular}{clcl}
\toprule
\textbf{Platform} & \textbf{Dataset Name / Description} & \textbf{Purpose} & \textbf{Notes} \\
\midrule
Roboflow & detect Dataset~\parencite{detect-lhuia_dataset} & Positive Class & Swastika and related Nazi symbols \\

Roboflow & swastika Dataset (imgann)~\parencite{swastika-o7bde_dataset} & Positive Class & Swastika symbol detection \\

Roboflow & swastika Dataset (Darja Nechaeva)~\parencite{swastika-uw8xt_dataset} & Positive Class & Historical swastika variations \\

Roboflow & swastika\_detect Dataset~\parencite{swastika_detect_dataset} & Positive Class & Swastika detection for moderation \\

Roboflow & symbols Dataset~\parencite{symbols-w93jw_dataset} & Positive Class & Includes general and hate symbols \\

Roboflow & violent object Dataset~\parencite{violent-object_dataset} & Positive Class & Hate symbols and violent imagery \\

Roboflow & Swastiki\_Net Dataset~\parencite{swastiki_net_dataset} & Positive Class & Swastika symbol classification \\

Roboflow & Content Moderation Dataset~\parencite{content-moderation-fwj5b_dataset} & Positive Class & Multiple harmful or extremist symbols \\

Roboflow & Swastika\_Detection Dataset~\parencite{swastika_detection_dataset} & Positive Class & Swastikas in different contexts \\

Kaggle & 100 Sports Image Classification~\parencite{100_sport_image_classification_dataset} & Negative Class & Neutral sports images \\

Kaggle & Car vs Bike Classification Dataset~\parencite{car_bike_classification_dataset} & Negative Class & Background class diversity \\

Kaggle & People Clothing Segmentation~\parencite{people_clothing_segmentation_dataset} & Negative Class & Background objects, neutral content \\

Kaggle & Fruits and Vegetables Dataset~\parencite{fruit_and_vegetable_image_recognition_dataset} & Negative Class & For reducing false positives \\

Kaggle & Animal Image Dataset~\parencite{animal_image_dataset} & Negative Class & Animal images for neutral comparison \\

Kaggle & Kids and Adults Detection~\parencite{kids_and_adults_detection_dataset} & Negative Class & Human figures, non-extremist \\

Kaggle & SCOOBY DOO Classification~\parencite{scooby_doo_classification_dataset} & Negative Class & Cartoon data for variance \\

Kaggle & Low Light Animals~\parencite{subrata_sarkar_2025} & Negative Class & Low-light conditions for robustness \\

Kaggle & AI vs Human-Generated Images~\parencite{ai_vs_human_generated_images_dataset} & Negative Class & Diverse image source simulation \\
\bottomrule
\end{tabular}
}
\caption{Summary of datasets used in this thesis for training, validation, and testing.}
\label{tab:dataset_summary}
\end{table}

\section{Nazi Symbol Description}
  \label{appendix:symbol-description}

The Nazi symbols are mainly from Roboflow with labels. The symbols are defined as follows:

\begin{itemize}
  \item \textbf{Sonnenrad}: The sonnenrad, or sunwheel, is an ancient symbol
  from Norse and Celtic cultures that was appropriated by the Nazis to
  promote an idealized "Aryan" heritage. Its many variations include forms
  of the swastika and Celtic Cross, and it was used by the Nazi Party,
  SA, and SS. Today, it is commonly used by neo-Nazis and white
  supremacists, especially a version with two concentric circles and
  crooked rays, sometimes featuring a swastika in the center.
  \item \textbf{Flash and circle}: The Flash and Circle is a symbol featuring
  a lightning bolt within a circle, originally used by the British Union
  of Fascists in the 1930s. It was meant to represent action within unity
  and has since been adopted by neo-fascist and white supremacist groups.
  Though less common today, it remains a recognizable hate symbol associated
  with fascist ideology.
  \item \textbf{Broken Sun Cross}: Symbols featuring an equilateral cross enclosed
  in a circle have appeared in various cultures since the Bronze Age. Since
  the 1930s, a version with ‘broken’ arms resembling a swastika gained far-right
  associations. It was used by the German Faith Movement and later as insignia
  for two Waffen SS divisions.
  \item \textbf{The Happy Merchant}: The Happy Merchant is an antisemitic meme
  depicting a Jewish man with exaggerated, stereotypical features, often used
  by white supremacists to promote hate. It originated from racist artwork
  by cartoonist Nick Bougas, who used the pseudonym “A. Wyatt Mann.” His
  offensive cartoons, widely circulated in the 1980s and 1990s, were often
  published by white supremacist Tom Metzger.
  \item \textbf{Hitler}: The images of Adolf Hilter.
  \item \textbf{Nazi salute}: The Nazi salute, also known as the Hitler salute
  or Sieg Heil salute, was a gesture used
  in Nazi Germany to show loyalty to Adolf Hitler and the Nazi Party, typically
  accompanied by phrases like "Heil Hitler" or "Sieg Heil". Inspired by the
  Italian Fascist salute, it became an official Nazi greeting in 1926 and
  was mandatory for civilians but optional for military personnel until
  1944. Today, the salute is banned in several countries—including Germany,
  Austria, and Slovakia—and is considered hate speech or a criminal act in
  many others if used to promote Nazi ideology.
%  \item \textbf{Neo-Nazi emblems}: Neo-Nazi emblems are symbols used by modern
%    white supremacist and neo-Nazi groups to express their extremist beliefs
%    and allegiance to Nazi ideology. These emblems often include modified or
%    obscure versions of historical Nazi symbols, such as the swastika, sonnenrad,
%    or runes, to avoid legal restrictions or public scrutiny. While some are
%    direct revivals, others are new creations that serve the same purpose of
%    spreading hate and promoting white nationalist ideas. In this study, it includes
%    symbols from various far-right extremist groups, such as Atomwaffen Division emblem,
%    Blood \& Honor emblem, Celtic Cross, Combat 18 emblem, Golden Dawn, Hammerskins emblem,
%    Identitarian movement emblem, Kolovrat, National Rebirth of Poland emblem,
%    Popular front emblem. Since only a limited number of images were available for
%    each group, they were consolidated under a single "Neo-Nazi" class for classification purposes.
  \item \textbf{Yellow badge}: The Yellow badge was a cloth patch, usually in
  the shape of a Star of David, that Jews were forced to wear by the Nazis
  to identify and segregate them. It symbolized persecution and was used to
  mark Jews publicly, contributing to their dehumanization and eventual
  deportation. Today, it stands as a powerful symbol of the Holocaust and
  the oppression of Jewish people.
  \item \textbf{Siegrune}: The Siegrune, or "S rune", is a lightning-bolt-shaped
  symbol used by the Nazi SS to represent victory ("Sieg" in German). It was
  adapted from ancient runic alphabets and became one of the most recognizable
  symbols of the SS, often displayed as a double rune on uniforms and insignia.
  Today, it is widely recognized as a hate symbol and is banned or restricted
  in several countries due to its association with Nazi ideology.
  \item \textbf{SS Skull}: The SS skull, also known as the "Totenkopf" ("death's head"),
  was a symbol used by the Nazi SS to represent loyalty unto death and intimidation.
  It was prominently featured on SS uniforms, especially those of the elite units
  and concentration camp guards. Today, it is considered a hate symbol due to
  its strong association with the crimes and ideology of the SS and the Nazi regime.
  \item \textbf{Sturmabteilung emblem}: The Sturmabteilung (SA) emblem was the
  symbol of the Nazi Party's paramilitary wing, known for its violent tactics
  and support for Adolf Hitler's rise to power. The emblem typically featured
  a stylized version of the SA initials, often accompanied by a swastika.
  It represented the SA's role in enforcing Nazi ideology and intimidation,
  and the emblem is now recognized as a symbol of hate and fascism.
  \item \textbf{Swastika}: The Swastika is an ancient symbol, originally used in
  various cultures to represent good luck, prosperity, and well-being. However,
  it was appropriated by the Nazi Party in the 1920s, becoming a powerful
  symbol of their ideology and representing white supremacy, hate, and
  anti-Semitism. Today, it is widely recognized as a symbol of Nazi oppression
  and is illegal in many countries due to its association with hate and violence.
  \item \textbf{Wolfsangel}: The Wolfsangel is a symbol originally used in medieval
  Europe, representing a wolf trap, and later adopted by Nazi and neo-Nazi
  groups. It was used by Nazi units during World War II, including the Waffen-SS,
  and has become associated with white supremacy and fascist ideologies. Today,
  the Wolfsangel is a recognized hate symbol, often used by extremist groups
  to promote Nazi beliefs.
  \item \textbf{Atomwaffen Division emblem}:The Atomwaffen Division emblem is a symbol
  used by a violent neo-Nazi terrorist group known for promoting accelerationism and
  engaging in acts of domestic terrorism. The emblem typically features a stylized
  radiation warning symbol, representing the group's fascination with chaos,
  destruction, and nuclear war. It is widely recognized as a hate symbol and is
  associated with extremist violence, white supremacy, and far-right terrorism.
  \item \textbf{Blood \& Honor emblem}: The Blood \& Honour emblem represents an international
  neo-Nazi network that promotes white supremacist music and ideology. The name comes
  from the Hitler Youth motto "Blut und Ehre" (Blood and Honor), and the emblem often
  includes those words along with fascist symbols like runes or laurel wreaths. It
  is widely recognized as a hate symbol linked to racist propaganda, violence,
  and far-right extremism.
  \item \textbf{Celtic Cross}: The Celtic Cross is an ancient symbol with religious
  and cultural origins in Celtic Christianity, featuring a cross with a circle
  around the intersection. While it remains a legitimate cultural and religious
  symbol, white supremacists have adopted a simplified version as a symbol of white
  identity and nationalism. This version is now widely recognized as a hate symbol
  in extremist contexts.
  \item \textbf{Combat 18 emblem}: Combat 18 is a violent neo-Nazi organization, with the
  number "18" representing the initials A and H (Adolf Hitler). Its emblem typically
  includes the number or stylized letters alongside other far-right imagery. The group
  is known for promoting white supremacy and terrorism, and its emblem is a symbol of
  hate and violence.
  \item \textbf{Golden Dawn}: Golden Dawn is a far-right political party from Greece known for its neo-Nazi ideology
  and involvement in violent, xenophobic acts. Its emblem resembles a Greek meander pattern, often compared to a
  swastika due to its stylized, angular design. Though once politically active, the group has been widely condemned
  and legally prosecuted for its extremist activities.
  \item \textbf{Hammerskins emblem}: The Hammerskins emblem features two crossed hammers, symbolizing the group's
  slogan "Hammerskins Nation" and white power. This violent white supremacist group is rooted in the skinhead
  movement and is heavily involved in hate music and racist activism. The emblem is a well-known hate symbol tied to
  global neo-Nazi networks.
  \item \textbf{Identitarian movement emblem}: The Identitarian Movement emblem typically features a yellow circle
  with a black Greek lambda symbol, inspired by the Spartan shield. This far-right European movement promotes
  anti-immigration and ethnonationalist ideologies under the guise of cultural preservation. Although presented with
  modern, clean branding, the emblem is associated with white nationalist beliefs.
  \item \textbf{Kolovrat}: The Kolovrat is an ancient Slavic symbol representing the sun and cosmic cycles, used
  historically in pagan contexts. In recent years, it has been appropriated by Slavic neo-Nazis and white
  supremacists as a symbol of ethnic purity and nationalist pride. While some use it in a cultural or historical
  sense, its association with extremist ideology marks it as a hate symbol in certain contexts.
  \item \textbf{National Rebirth of Poland emblem}: The National Rebirth of Poland (NOP) emblem represents a
  far-right nationalist party in Poland known for its anti-Semitic and xenophobic views. The emblem often
  incorporates stylized nationalist imagery, including swords or eagles, designed to evoke Polish strength and
  purity. It is associated with extremist rhetoric and far-right political activism in Poland.
  \item \textbf{Popular front emblem}: The Popular Front emblem is used by a range of nationalist or fascist groups,
  often combining traditional nationalist symbols like eagles or stars with militaristic imagery. Though the term "
  Popular Front" historically referred to leftist coalitions, in extremist contexts, it can represent
  ultra-nationalist, anti-democratic movements. The emblem typically signals far-right or authoritarian ideologies.
%  \item \textbf{Doppelsiegrune}: The Doppelsiegrune, or "double Siegrune," features two lightning-bolt-like runes
%  and was the official insignia of the Nazi SS. It symbolized power, obedience, and loyalty to Adolf Hitler and
%  became one of the most infamous emblems of the Nazi regime. Today, it is banned or heavily restricted in many
%  countries and is recognized worldwide as a symbol of hate and terror.
\end{itemize}